\section{Related Work}
\label{sec:related}

\para{Hidden-state prediction of future tokens.}
\citet{pal2023future} introduced \futurelens, demonstrating that a single hidden state in GPT-J-6B can predict tokens two or more positions ahead with up to 48.4\% accuracy using learned prompt causal interventions.
Prediction accuracy peaks at middle layers for future tokens (unlike next-token prediction, which peaks at the final layer), suggesting that future information is pre-computed during forward passes.
Our work extends this finding by testing whether the strength of this pre-computation correlates with non-compositionality.

\para{Token erasure and implicit vocabularies.}
\citet{feucht2024token} discovered that the last token position of multi-token words and named entities exhibits ``erasure''---information about preceding tokens is rapidly forgotten in early layers as the model overwrites token-level representations with entity-level semantics.
Their erasure score identifies multi-token lexical items from model internals.
Where \futurelens detects forward prediction (can the model anticipate tokens ahead?), \tokenerasure detects backward forgetting (does the model discard constituent information?).
Both signals indicate that the model treats a phrase as a single unit, and combining them is a promising direction we discuss in \secref{sec:discussion}.

\para{Memorization recall in transformers.}
\citet{haviv2022understanding} provided the first framework for probing memorization recall using idioms.
They constructed the \idiomem dataset of 852 English idioms and discovered a two-phase memorization profile: early-layer \emph{candidate promotion} (the target token rises to the top of the distribution) followed by late-layer \emph{confidence boosting} (probability sharply increases).
We use \idiomem as a validation set of known \noncomp phrases and extend the analysis with \futurelens metrics beyond the \logitlens used in their work.

\para{Multi-word expression semantics in transformers.}
\citet{miletic2024systematic} surveyed how transformers capture MWE semantics, finding that models rely inconsistently on surface patterns and memorized information, with MWE meaning predominantly localized in early layers.
\citet{rambelli2023frequent} showed that idioms and high-frequency \comp phrases are processed similarly by both humans and neural language models (measured via surprisal), raising an important confound: frequency-based retrieval may mimic non-compositionality signals.
We address this concern by using expert-annotated compositionality labels from the \farahmand dataset, which contains compounds spanning a range of frequencies.

\para{Compositionality annotation.}
\citet{farahmand2015multiword} annotated 1,042 English noun compounds with binary compositionality judgments from four expert annotators, providing the controlled evaluation set we use.
\citet{reddy2011empirical} provided continuous compositionality scores for 90 noun compounds, offering finer granularity but smaller scale.
We use the Farahmand dataset for its larger size and the availability of majority-vote labels for binary classification.

\para{Logit lens interpretation.}
The \logitlens~\citep{nostalgebraist2020logitlens} applies the model's output head (layer normalization followed by the unembedding matrix) to intermediate hidden states, interpreting them as probability distributions over the vocabulary.
\citet{din2023jump} and \citet{geva2022transformer} used this approach to trace how predictions evolve across layers.
We use the \logitlens as one of our three probing signals, applying it at nine layers spanning the full depth of \gptxl.
